% =============================================================
%  ch06_schrittmotoren.tex  --  Schrittmotoren und Drehmomentschätzung
% =============================================================

\chapter{Schrittmotoren und Drehmomentschätzung}

\section{NEMA-17-Grundlagen}
Ein NEMA-17 ist ein bipolarer Hybridschrittmotor mit typisch
\textbf{1{,}8\textdegree/Schritt} = 200 Vollschritte/Umdrehung. Per
\textbf{Microstepping} (z.\,B.\ 1/16, 1/256 beim TMC) wird sehr feinaufgelöst
und vibrationsarm gefahren. Phasenströme typisch 0{,}4--2\,A, Haltemoment je
nach Baulänge $\approx$ 0{,}2--0{,}6\,Nm. Im \textbf{offenen Regelkreis} kennt
der Motor seine Last nicht -- bei Überlast verliert er Schritte unbemerkt.

\section{Indirekte Drehmomentschätzung}
Statt eines teuren Drehmomentsensors wird das Lastmoment aus den
\emph{elektrischen} Größen geschätzt. Physikalischer Hintergrund: Bei einem
belasteten Schrittmotor wächst der \textbf{Lastwinkel} (die Verschiebung
zwischen Statorfeld und Rotor). Dieser zeigt sich in Phasenströmen und
Gegen-EMK. Zwei praktische Umsetzungen:
\begin{enumerate}
\item \textbf{StallGuard (Trinamic):} Der Treiber misst intern die Gegen-EMK und
  liefert einen lastproportionalen Zahlenwert (\code{SG\_RESULT}) über UART.
  Das ist der bevorzugte Ansatz -- ein Proxy fürs Drehmoment ohne
  Zusatzhardware.
\item \textbf{Externe Strom-/Spannungsmessung:} Ein Sensormodul (z.\,B.\
  INA226) misst Busspannung und -strom pro Treiber. Zusammen mit dem bekannten
  Microstep-Zustand lässt sich die mechanische Leistung und damit das Moment
  abschätzen.
\end{enumerate}
Beide Wege sind weniger präzise als ein echter Sensor, aber günstig und
praktikabel -- siehe Teil~II für die konkrete Hardware.
